\section{Conclusion}\label{sec:conclusion}

We presented a post-dissertation methodological refinement and
generalization of the DoE + RSM + Varimax + NBI framework for
multi-objective hyperparameter tuning of GBDT models. The article-track
implementation departs from the dissertation pipeline of
\citep{ribeiro2026dissertation} in four substantive ways:
(i) the optimizer is a true $N$-objective NBI rather than a
weighted-sum scalarization; (ii) every objective carries an explicit
direction, with FMSE wrapping for $\VRF$ targets; (iii) the surrogate
fitting is design-aware, separating process-quadratic RSM from
Scheff\'{e} mixture polynomials; and (iv) a conditional MBPA
post-optimization stage refines weights when the primary frontier is
weakly informative. The legacy weighted-sum solver is retained as an
ablation baseline and is no longer named NBI in either the code base
or the article text.

\todoblock{Quantitative claims (CPU-hour reduction, predictive parity
intervals, MBPA trigger rate, NBI residual percentiles) go here once
the v1 campaign is aggregated. Avoid hyperbole; report medians and
inter-quartile ranges, not single-replica winners.}

\paragraph{Future work.} Three directions follow naturally:
the doctoral 82-dataset benchmark; an extension to non-tree
estimators with the same \texttt{ObjectiveSpec} layer
(linear models, deep tabular networks); and a tighter coupling
between the cost estimator and the orchestrator so that a campaign
can self-tune its replica count to a wall-clock budget.
